\lecture{04}{Основные постановки задач оптимизации}

\sourcevideo
    {https://www.youtube.com/playlist?list=PLOzRYVm0a65fhKDS8cYIC7YCuVGJ4FOJx}
    {Week 1: Lecture 4: Basics of Optimization Problems}

Перед построением распределённых алгоритмов необходимо определить,
какая задача решается, какие величины являются неизвестными и какие
ограничения должны выполняться. В этой лекции вводится общая
математическая постановка оптимизации и рассматриваются две
прикладные модели: разреженное восстановление сигнала и метод
опорных векторов.

\section{Общая задача оптимизации и ограничения}

Пусть
\[
    f\colon\RR^n\to\RR.
\]
Задача без ограничений имеет вид
\[
    \min_{x\in\RR^n} f(x).
\]

\begin{definition}[Целевая функция]
Функция \(f\), значение которой требуется минимизировать или
максимизировать, называется целевой функцией.
\end{definition}

\begin{definition}[Переменная решения]
Вектор
\[
    x=
    \begin{pmatrix}
        x_1&\cdots&x_n
    \end{pmatrix}^{\mathsf T}
    \in\RR^n
\]
называется переменной решения, оптимизационной переменной или
вектором параметров задачи.
\end{definition}

\begin{definition}[Задача без ограничений]
Задача
\[
    \min_{x\in\RR^n}f(x)
\]
называется задачей без ограничений, если переменная \(x\) может
принимать любое значение из области определения \(f\).
\end{definition}

Максимизацию можно свести к минимизации:
\[
    \max_x f(x)
    \quad\Longleftrightarrow\quad
    \min_x\bigl(-f(x)\bigr).
\]
Поэтому далее достаточно рассматривать задачи минимизации.



Общая задача с ограничениями записывается как
\[
    \begin{aligned}
        \min_{x\in\RR^n}
        \quad&
        f(x),
        \\
        \text{при условиях}
        \quad&
        h_i(x)\le0,
        \qquad i=1,\ldots,m,
        \\
        &
        \ell_j(x)=0,
        \qquad j=1,\ldots,r.
    \end{aligned}
\]

Функции \(h_i\) задают ограничения-неравенства, а функции
\(\ell_j\) --- ограничения-равенства.

\begin{definition}[Допустимое множество]
Множество
\[
    \mathcal X
    =
    \left\{
        x\in\RR^n:
        \begin{array}{l}
            h_i(x)\le0,\quad i=1,\ldots,m,\\
            \ell_j(x)=0,\quad j=1,\ldots,r
        \end{array}
    \right\}
\]
называется допустимым множеством задачи.
\end{definition}

\begin{definition}[Допустимая точка]
Точка \(x\in\mathcal X\), удовлетворяющая всем ограничениям,
называется допустимой.
\end{definition}

\begin{definition}[Оптимальное решение]
Допустимая точка \(x^\star\in\mathcal X\) называется глобальным
решением задачи, если
\[
    f(x^\star)\le f(x)
    \qquad
    \text{для всех }x\in\mathcal X.
\]
Соответствующее значение
\[
    f^\star=f(x^\star)
\]
называется оптимальным значением.
\end{definition}

Таким образом, задача оптимизации состоит не в поиске произвольной
допустимой точки, а в выборе среди допустимых точек той, на которой
целевая функция принимает наименьшее значение.

\section{Выпуклая задача и методы оптимизации}

Задача
\[
    \begin{aligned}
        \min_x
        \quad& f(x),
        \\
        \text{при условиях}
        \quad& h_i(x)\le0,
        \\
        & \ell_j(x)=0
    \end{aligned}
\]
является выпуклой, если \(f\) и все функции \(h_i\) выпуклы, а
ограничения-равенства аффинны:
\[
    \ell_j(x)=a_j^{\mathsf T}x-b_j.
\]

В этом случае допустимое множество выпукло, а всякий локальный
минимум является глобальным. Формальные определения выпуклых
множеств и функций вводятся в следующей лекции.



Оптимизационные методы различаются по информации о целевой функции,
которую они используют.

\begin{definition}[Метод первого порядка]
Методом первого порядка называется алгоритм, использующий значения
функции и её градиента
\[
    \nabla f(x).
\]
\end{definition}

Типичным примером является градиентный спуск:
\[
    x^{k+1}=x^k-\eta_k\nabla f(x^k).
\]

\begin{definition}[Метод второго порядка]
Методом второго порядка называется алгоритм, дополнительно
использующий матрицу Гессе
\[
    \nabla^2f(x).
\]
\end{definition}

Если матрица Гессе в точке \(x^k\) обратима, классический шаг
метода Ньютона имеет вид
\[
    x^{k+1}
    =
    x^k
    -
    \bigl(\nabla^2f(x^k)\bigr)^{-1}
    \nabla f(x^k).
\]

Методы второго порядка учитывают локальную кривизну функции, но
вычисление, хранение и обращение матрицы Гессе может быть дорого при
большой размерности. К методам, использующим приближение информации
второго порядка, относятся квазиньютоновские методы, в частности
L-BFGS. В лекции также упоминается AdaHessian.

\section{Разреженное восстановление сигнала}

Пусть неизвестный сигнал \(x\in\RR^n\) наблюдается через линейную
модель
\[
    y=Mx+v,
\]
где
\[
    y\in\RR^s,\qquad
    M\in\RR^{s\times n},\qquad
    v\in\RR^s.
\]
Здесь \(y\) --- измеренный сигнал, \(M\) --- известная матрица
измерений, а \(v\) --- шум.

Если \(s\ll n\), линейная система недоопределена: число измерений
меньше числа неизвестных. Для восстановления \(x\) требуется
дополнительное структурное предположение. В задачах разреженного
восстановления предполагается, что большинство компонент \(x\)
равно нулю.

\begin{definition}[Разреженный вектор]
Носителем вектора \(x\in\RR^n\) называется множество
\[
    \operatorname{supp}(x)
    \coloneqq
    \{i\in\{1,\ldots,n\}:x_i\ne0\}.
\]
Вектор называется разреженным, если мощность его носителя мала по
сравнению с \(n\). Число ненулевых компонент обозначают через
\[
    \lVert x\rVert_0
    \coloneqq
    \bigl\lvert\operatorname{supp}(x)\bigr\rvert,
\]
хотя \(\lVert\cdot\rVert_0\) не является нормой.
\end{definition}

Оптимизация с непосредственным штрафом \(\lVert x\rVert_0\) в общем
случае невыпукла и комбинаторно сложна. Стандартная выпуклая замена ---
\(\ell_1\)-регуляризация.

\begin{problem}[\(\ell_1\)-регуляризованный метод наименьших квадратов]
Требуется решить задачу
\[
    \min_{x\in\RR^n}
    \left\{
        \frac12\lVert Mx-y\rVert_2^2
        +
        \lambda\lVert x\rVert_1
    \right\},
    \qquad
    \lambda>0.
\]
\end{problem}

Первое слагаемое измеряет согласованность восстановленного сигнала с
наблюдениями. Второе слагаемое
\[
    \lVert x\rVert_1
    =
    \sum_{j=1}^{n}\lvert x_j\rvert
\]
штрафует суммарную величину компонент и тем самым способствует
разреженности решения. Параметр \(\lambda\) задаёт компромисс между
точностью восстановления и силой регуляризации.

Задача не содержит явных ограничений и потому относится к задачам
без ограничений. При этом слагаемое \(\lVert x\rVert_1\) негладко в
точках с нулевыми компонентами, поэтому применяются, например,
субградиентные или проксимальные методы. Такие модели используются в
сжатых измерениях, обработке изображений и разреживании параметров
машинного обучения.



Пусть один и тот же неизвестный сигнал \(x\) наблюдается \(N\)
агентами:
\[
    y_i=M_ix+v_i,
    \qquad
    i=1,\ldots,N.
\]
Агент \(i\) знает только \(M_i\) и \(y_i\). Централизованная задача
имеет вид
\[
    \min_{x\in\RR^n}
    \left\{
        \frac12
        \sum_{i=1}^{N}
        \lVert M_ix-y_i\rVert_2^2
        +
        \lambda\lVert x\rVert_1
    \right\}.
\]

Для распределённой записи каждому агенту сопоставляется локальная
копия \(x_i\):
\[
    \begin{aligned}
        \min_{x_1,\ldots,x_N}
        \quad&
        \sum_{i=1}^{N}
        \left(
            \frac12\lVert M_ix_i-y_i\rVert_2^2
            +
            \frac{\lambda}{N}\lVert x_i\rVert_1
        \right),
        \\
        \text{при условиях}
        \quad&
        x_1=x_2=\cdots=x_N.
    \end{aligned}
\]

\begin{remark}
Множитель \(\lambda/N\) распределяет исходный регуляризатор между
агентами. При выполнении условия консенсуса \(x_i=x\) сумма локальных
регуляризаторов равна \(\lambda\lVert x\rVert_1\).
\end{remark}

Исходная централизованная задача без ограничений после введения
локальных копий становится задачей с ограничениями. Эти ограничения
необходимы, поскольку все агенты восстанавливают один и тот же
физический сигнал.

Если каждый агент решает только собственную задачу
\[
    \min_{x_i}
    \left\{
        \frac12\lVert M_ix_i-y_i\rVert_2^2
        +
        \frac{\lambda}{N}\lVert x_i\rVert_1
    \right\},
\]
то локальные оценки \(x_i^\star\) в общем случае различаются.
Коммуникации позволяют учитывать измерения других узлов без передачи
всего локального набора данных.

Для связного неориентированного графа условие глобального консенсуса
эквивалентно локальным равенствам
\[
    x_i=x_j,
    \qquad
    (i,j)\in\Edges.
\]
Связность графа тогда влечёт
\[
    x_1=\cdots=x_N.
\]

Распределённый алгоритм должен одновременно уменьшать локальные
невязки
\[
    \lVert M_ix_i-y_i\rVert_2
\]
и согласовывать оценки соседних агентов. Поэтому распределённая
оптимизация объединяет два процесса:
\[
    \text{оптимизацию локальных функций}
    \quad\text{и}\quad
    \text{достижение консенсуса}.
\]

\section{Линейная классификация и геометрия гиперплоскости}

Рассмотрим обучающую выборку
\[
    \{(z_i,y_i)\}_{i=1}^{m},
    \qquad
    z_i\in\RR^n,
    \qquad
    y_i\in\{-1,+1\}.
\]
Здесь \(z_i\) обозначает вектор признаков, а \(y_i\) --- метку класса.

Линейный классификатор задаётся гиперплоскостью
\[
    w^{\mathsf T}z+b=0,
\]
где
\[
    w\in\RR^n,\qquad b\in\RR.
\]
Вектор \(w\) является нормалью к гиперплоскости, а \(b\) задаёт её
смещение относительно начала координат.

Предсказанная метка определяется знаком:
\[
    \widehat y(z)
    =
    \operatorname{sign}\bigl(w^{\mathsf T}z+b\bigr).
\]
Для правильной классификации обучающей точки требуется
\[
    y_i\bigl(w^{\mathsf T}z_i+b\bigr)>0.
\]

Через одни и те же линейно разделимые данные можно провести несколько
разделяющих гиперплоскостей. Метод опорных векторов выбирает ту, для
которой минимальное расстояние от обучающих точек до гиперплоскости
максимально.



При \(w\ne0\) расстояние от точки \(z_i\) до гиперплоскости
\[
    w^{\mathsf T}z+b=0
\]
равно
\[
    d_i
    =
    \frac{
        \lvert w^{\mathsf T}z_i+b\rvert
    }{
        \lVert w\rVert_2
    }.
\]

Для правильно классифицированной точки
\[
    \lvert w^{\mathsf T}z_i+b\rvert
    =
    y_i\bigl(w^{\mathsf T}z_i+b\bigr),
\]
поэтому минимальный геометрический отступ равен
\[
    \gamma
    =
    \min_{1\le i\le m}
    \frac{
        y_i(w^{\mathsf T}z_i+b)
    }{
        \lVert w\rVert_2
    }.
\]

Параметры гиперплоскости определены неоднозначно: для любого
\(\alpha>0\) пары
\[
    (w,b)
    \qquad\text{и}\qquad
    (\alpha w,\alpha b)
\]
задают одну и ту же гиперплоскость. Если выборка линейно разделима,
положительным масштабированием параметры нормируют так, чтобы
\[
    \min_{1\le i\le m}
    y_i(w^{\mathsf T}z_i+b)
    =
    1.
\]
Тогда ограничения принимают вид
\[
    y_i(w^{\mathsf T}z_i+b)\ge1,
    \qquad i=1,\ldots,m,
\]
а геометрический отступ равен
\[
    \gamma=\frac{1}{\lVert w\rVert_2}.
\]

Максимизация \(\gamma\) эквивалентна минимизации нормы \(w\):
\[
    \max_{w,b}
    \frac{1}{\lVert w\rVert_2}
    \quad\Longleftrightarrow\quad
    \min_{w,b}
    \frac12\lVert w\rVert_2^2.
\]

\section{Метод опорных векторов и распределённая классификация}

\begin{problem}[SVM с жёстким зазором (hard margin)]
Для линейно разделимой выборки требуется решить задачу
\[
    \begin{aligned}
        \min_{w\in\RR^n,\ b\in\RR}
        \quad&
        \frac12\lVert w\rVert_2^2,
        \\
        \text{при условиях}
        \quad&
        y_i(w^{\mathsf T}z_i+b)\ge1,
        \qquad i=1,\ldots,m.
    \end{aligned}
\]
\end{problem}

Минимизация целевой функции при заданной нормировке эквивалентна
максимизации геометрического отступа, а ограничения требуют правильной
классификации всех обучающих точек.

\begin{definition}[Опорный вектор]
Обучающая точка \(z_i\) называется опорным вектором оптимального
классификатора, если соответствующее ограничение активно:
\[
    y_i\bigl(({w^\star})^{\mathsf T}z_i+b^\star\bigr)=1.
\]
\end{definition}

Опорные векторы являются ближайшими к разделяющей гиперплоскости
точками. Именно они определяют положение гиперплоскости максимального
отступа.

\begin{remark}
Функция
\[
    \frac12\lVert w\rVert_2^2
\]
квадратична и сильно выпукла по \(w\), но не зависит от \(b\) и
потому не является сильно выпуклой по паре \((w,b)\). Такая форма
удобнее для анализа и вычислений, чем непосредственная максимизация
\(1/\lVert w\rVert_2\).
\end{remark}



Пусть обучающие данные распределены между \(N\) агентами:
\[
    \Dataset_i
    =
    \{(z_{ij},y_{ij})\}_{j=1}^{m_i}.
\]
Агент \(i\) видит только собственную выборку \(\Dataset_i\).

Если каждый агент независимо строит локальный классификатор, он
решает
\[
    \begin{aligned}
        \min_{w_i,b_i}
        \quad&
        \frac12\lVert w_i\rVert_2^2,
        \\
        \text{при условиях}
        \quad&
        y_{ij}(w_i^{\mathsf T}z_{ij}+b_i)\ge1,
        \qquad j=1,\ldots,m_i.
    \end{aligned}
\]
Полученные гиперплоскости могут существенно различаться, поскольку
локальные выборки содержат разные участки общего распределения
данных.

Централизованный классификатор должен учитывать объединённую выборку
\[
    \Dataset
    =
    \bigcup_{i=1}^{N}\Dataset_i.
\]
Если объединённая выборка линейно разделима, её распределённая
постановка имеет вид
\[
    \begin{aligned}
        \min_{\substack{w_1,\ldots,w_N\\b_1,\ldots,b_N}}
        \quad&
        \frac{1}{2N}
        \sum_{i=1}^{N}\lVert w_i\rVert_2^2,
        \\
        \text{при условиях}
        \quad&
        y_{ij}(w_i^{\mathsf T}z_{ij}+b_i)\ge1,
        \qquad
        j=1,\ldots,m_i,
        \\
        &
        w_i=w_\ell,\qquad
        b_i=b_\ell,
        \qquad
        (i,\ell)\in\Edges.
    \end{aligned}
\]

При связном неориентированном графе локальные равенства приводят к
общим параметрам:
\[
    w_1=\cdots=w_N,
    \qquad
    b_1=\cdots=b_N.
\]
После подстановки консенсуса коэффициент \(1/N\) сокращает сумму
одинаковых квадратичных слагаемых, поэтому распределённая задача
эквивалентна централизованной. При этом каждый агент использует только
собственные данные и сообщения соседей.

Практическая сложность состоит в том, что локально оптимальная
гиперплоскость может плохо классифицировать данные, отсутствующие у
данного агента. Поэтому обмен информацией необходим не только для
согласования параметров, но и для приближения решения, построенного
по совокупной выборке.

\section{Неразделимые данные}

Жёстко-разделяющая постановка требует существования \(w,b\), для
которых
\[
    y_i(w^{\mathsf T}z_i+b)\ge1
\]
выполняется для всех точек. Если классы линейно неразделимы или
некоторые метки ошибочны, допустимое множество такой задачи может
оказаться пустым.

В лекции отмечаются два стандартных обобщения. SVM с мягким зазором
допускает нарушения ограничений и штрафует их величину. Ядерный SVM
неявно работает в пространстве признаков, где разделяющая поверхность
может быть линейной. Подробный вывод этих моделей в лекции не
проводится.


Общая задача оптимизации задаётся целевой функцией, переменной решения
и системой ограничений. Допустимое множество содержит точки,
удовлетворяющие всем ограничениям. Методы первого порядка используют
градиенты, а методы второго порядка дополнительно используют
кривизну функции. В разреженном восстановлении \(\ell_1\)-регуляризация
способствует получению решений с небольшим числом ненулевых компонент. При
распределении измерений между агентами вводятся локальные копии
неизвестного сигнала и ограничения консенсуса. В методе опорных
векторов максимизация минимального геометрического отступа приводит
к выпуклой квадратичной задаче
\[
    \begin{aligned}
        \min_{w,b}\quad&\frac12\lVert w\rVert_2^2,\\
        \text{при условиях}\quad&
        y_i(w^{\mathsf T}z_i+b)\ge1,
        \qquad i=1,\ldots,m.
    \end{aligned}
\]
При распределённых данных локальные классификаторы могут существенно
отличаться от решения по объединённой выборке, поэтому оптимизация
должна сочетаться с коммуникацией и согласованием параметров.
